% GENERATED FILE. Edit canonical lessons, then run npm run generate:books.
% canonical-source-hash: fnv1a64:af6bda34f2a299d7
% canonical-lessons: MR-C433-orakhda, MR-C433-tada, MR-C433-ganjne, MR-C433-cepne, MR-C433-band-padne

\chapter{Rust and Dents}
\label{ch:mr-rust-and-dents}
\hlchaptermodality{433}

\begin{chapteropening}
\textbf{By the end of this chapter:} \emph{I can talk about scratches and cracks, and say a thing rusted, got dented or broke down.}

Complete the last lesson of chapter 433: I can talk about scratches and cracks, and say a thing rusted, got dented or broke down.
\end{chapteropening}

\section[\texorpdfstring{\mr{ओरखडा} (orakhḍā)}{ओरखडा (orakhḍā)}]{\mr{ओरखडा} (orakhḍā) — a scrape; a scratch}

\label{lesson:MR-C433-orakhda}



\begin{quote}
Before the new one: say the Marathi for to get messed up, then the Marathi for to ooze.
\end{quote}

\subsection*{You'll want to know: \mr{ओरखडा}}
\textbf{\mr{ओरखडा}} — \emph{orakhḍā} — \textquotedblleft{}a scrape; a scratch\textquotedblright{}.

\begin{quote}
\mr{गाडीच्या} \mr{दारावर} \mr{मोठा} \mr{ओरखडा} \mr{आहे}.
\end{quote}

\begin{quote}
\emph{gāḍīcyā dārāvar moṭhā orakhḍā āhe.}
\end{quote}

\begin{quote}
There is a big scratch on the car door.
\end{quote}

Masculine.

\mr{ओरखडा} \mr{पडणे}, a scratch appears.

\begin{cousinweb}
\textbf{\mr{ओरखडा}} (\emph{orakhḍā}) is a Marathi word for a long, thin scratch; its origin is not settled.
\end{cousinweb}

\begin{grammarlens}[title={What you've built}]
One more word for getting around.
\end{grammarlens}

\subsection*{Your turn}
\begin{itemize}
\raggedright
  \item \emph{Say it:} \emph{orakhḍā}
  \item \emph{Say it:} \emph{orakhḍā}, once more
  \item \emph{Say it:} \emph{gāḍīcyā dārāvar moṭhā orakhḍā āhe.} — There is a big scratch on the car door.
  \item \emph{Recall:} say the Marathi for to get messed up, then the Marathi for to ooze, then say \emph{orakhḍā} again
\end{itemize}

\begin{tcolorbox}[breakable,colback=teal!4,colframe=teal!35!black,title={Before you move on}]
What does \mr{ओरखडा} mean? (\textquotedblleft{}A scrape; a scratch\textquotedblright{}.) Say it once more.
\end{tcolorbox}
\section[\texorpdfstring{\mr{तडा} (taḍā)}{तडा (taḍā)}]{\mr{तडा} (taḍā) — a split; a crack}

\label{lesson:MR-C433-tada}



\begin{quote}
Before the new one: say the Marathi for to ooze, then the Marathi for a scrape.
\end{quote}

\subsection*{You'll want to know: \mr{तडा}}
\textbf{\mr{तडा}} — \emph{taḍā} — \textquotedblleft{}a split; a crack\textquotedblright{}.

\begin{quote}
\mr{भूकंपानंतर} \mr{भिंतीला} \mr{तडा} \mr{गेला}.
\end{quote}

\begin{quote}
\emph{bhūkampānantar bhintīlā taḍā gelā.}
\end{quote}

\begin{quote}
After the earthquake the wall cracked.
\end{quote}

Masculine.

\mr{तडा} \mr{जाणे}, to crack: \mr{भिंतीला} \mr{तडा} \mr{गेला}.

\begin{cousinweb}
\textbf{\mr{तडा}} (\emph{taḍā}) is inherited: CDIAL lists Marathi \emph{taḍā}, \textquotedblleft{}a fissure\textquotedblright{}, beside \emph{taḍṇẽ}, \textquotedblleft{}to crack, split\textquotedblright{}, under \emph{*traṭ}, \textquotedblleft{}crackle\textquotedblright{} (CDIAL 5988).
\end{cousinweb}

\begin{grammarlens}[title={What you've built}]
One more word for getting around.
\end{grammarlens}

\subsection*{Your turn}
\begin{itemize}
\raggedright
  \item \emph{Say it:} \emph{taḍā}
  \item \emph{Say it:} \emph{taḍā}, once more
  \item \emph{Say it:} \emph{bhūkampānantar bhintīlā taḍā gelā.} — After the earthquake the wall cracked.
  \item \emph{Recall:} say the Marathi for to ooze, then the Marathi for a scrape, then say \emph{taḍā} again
\end{itemize}

\begin{tcolorbox}[breakable,colback=teal!4,colframe=teal!35!black,title={Before you move on}]
What does \mr{तडा} mean? (\textquotedblleft{}A split; a crack\textquotedblright{}.) Say it once more.
\end{tcolorbox}
\section[\texorpdfstring{\mr{गंजणे} (gañjṇe)}{गंजणे (gañjṇe)}]{\mr{गंजणे} (gañjṇe) — to rust}

\label{lesson:MR-C433-ganjne}



\begin{quote}
Before the new one: say the Marathi for a scrape, then the Marathi for a split.
\end{quote}

\subsection*{You'll want to know: \mr{गंजणे}}
\textbf{\mr{गंजणे}} — \emph{gañjṇe} — \textquotedblleft{}to rust\textquotedblright{}.

\begin{quote}
\mr{पावसात} \mr{सायकल} \mr{गंजली}.
\end{quote}

\begin{quote}
\emph{pāvsāt sāykal gañjlī.}
\end{quote}

\begin{quote}
The bicycle rusted in the rain.
\end{quote}

Past (he …): \mr{तो} \mr{गंजला}, \emph{to gañjlā}.

The noun is \mr{गंज}, rust: \mr{गंज} \mr{चढणे}, rust sets in.

\begin{cousinweb}
\textbf{\mr{गंजणे}} (\emph{gañjṇe}) is built on \textbf{\mr{गंज}}, \textquotedblleft{}rust\textquotedblright{}, whose origin is not settled.
\end{cousinweb}

\begin{grammarlens}[title={What you've built}]
One more word for getting around.
\end{grammarlens}

\subsection*{Your turn}
\begin{itemize}
\raggedright
  \item \emph{Say it:} \emph{gañjṇe}
  \item \emph{Say it:} \emph{gañjṇe}, once more
  \item \emph{Say it:} \emph{pāvsāt sāykal gañjlī.} — The bicycle rusted in the rain.
  \item \emph{Recall:} say the Marathi for a scrape, then the Marathi for a split, then say \emph{gañjṇe} again
\end{itemize}

\begin{tcolorbox}[breakable,colback=teal!4,colframe=teal!35!black,title={Before you move on}]
What does \mr{गंजणे} mean? (\textquotedblleft{}To rust\textquotedblright{}.) Say it once more.
\end{tcolorbox}
\section[\texorpdfstring{\mr{चेपणे} (cepṇe)}{चेपणे (cepṇe)}]{\mr{चेपणे} (cepṇe) — to get dented}

\label{lesson:MR-C433-cepne}



\begin{quote}
Before the new one: say the Marathi for a split, then the Marathi for to rust.
\end{quote}

\subsection*{You'll want to know: \mr{चेपणे}}
\textbf{\mr{चेपणे}} — \emph{cepṇe} — \textquotedblleft{}to get dented\textquotedblright{}.

\begin{quote}
\mr{भिंतीला} \mr{आपटून} \mr{गाडीचा} \mr{दरवाजा} \mr{चेपला}.
\end{quote}

\begin{quote}
\emph{bhintīlā āpṭūn gāḍīcā darvājā ceplā.}
\end{quote}

\begin{quote}
The car door got dented when it hit the wall.
\end{quote}

Past (he …): \mr{तो} \mr{चेपला}, \emph{to ceplā}.

It also means to press or squeeze something.

\begin{cousinweb}
\textbf{\mr{चेपणे}} (\emph{cepṇe}) is inherited: CDIAL lists Marathi \emph{cepṇẽ}, \textquotedblleft{}to press, squeeze\textquotedblright{}, under \emph{*capp}, \textquotedblleft{}press\textquotedblright{} (CDIAL 4674).
\end{cousinweb}

\begin{grammarlens}[title={What you've built}]
One more word for getting around.
\end{grammarlens}

\subsection*{Your turn}
\begin{itemize}
\raggedright
  \item \emph{Say it:} \emph{cepṇe}
  \item \emph{Say it:} \emph{cepṇe}, once more
  \item \emph{Say it:} \emph{bhintīlā āpṭūn gāḍīcā darvājā ceplā.} — The car door got dented when it hit the wall.
  \item \emph{Recall:} say the Marathi for a split, then the Marathi for to rust, then say \emph{cepṇe} again
\end{itemize}

\begin{tcolorbox}[breakable,colback=teal!4,colframe=teal!35!black,title={Before you move on}]
What does \mr{चेपणे} mean? (\textquotedblleft{}To get dented\textquotedblright{}.) Say it once more.
\end{tcolorbox}
\section[\texorpdfstring{\mr{बंद} \mr{पडणे} (band paḍṇe)}{बंद पडणे (band paḍṇe)}]{\mr{बंद} \mr{पडणे} (band paḍṇe) — to stop working; to break down}

\label{lesson:MR-C433-band-padne}



\begin{quote}
Before the new one: say the Marathi for to rust, then the Marathi for to get dented.
\end{quote}

\subsection*{You'll want to know: \mr{बंद} \mr{पडणे}}
\textbf{\mr{बंद} \mr{पडणे}} — \emph{band paḍṇe} — \textquotedblleft{}to stop working; to break down\textquotedblright{}.

\begin{quote}
\mr{घाटात} \mr{आमची} \mr{बस} \mr{बंद} \mr{पडली}.
\end{quote}

\begin{quote}
\emph{ghāṭāt āmcī bas band paḍlī.}
\end{quote}

\begin{quote}
Our bus broke down on the ghat road.
\end{quote}

Past (he …): \mr{तो} \mr{बंद} \mr{पडला}, \emph{to band paḍlā}.

Machines, clocks and lifts \mr{बंद} \mr{पडतात}; \mr{बंद} \mr{करणे} is to switch off.

\begin{cousinweb}
\textbf{\mr{बंद} \mr{पडणे}} (\emph{band paḍṇe}) joins \textbf{\mr{बंद}}, \textquotedblleft{}closed\textquotedblright{}, a loan usually traced to Persian \emph{band}, to \textbf{\mr{पडणे}}, \textquotedblleft{}to fall\textquotedblright{}: to fall still.
\end{cousinweb}

\begin{grammarlens}[title={What you've built}]
One more word for getting around.
\end{grammarlens}

\subsection*{Your turn}
\begin{itemize}
\raggedright
  \item \emph{Say it:} \emph{band paḍṇe}
  \item \emph{Say it:} \emph{band paḍṇe}, once more
  \item \emph{Say it:} \emph{ghāṭāt āmcī bas band paḍlī.} — Our bus broke down on the ghat road.
  \item \emph{Recall:} say the Marathi for to rust, then the Marathi for to get dented, then say \emph{band paḍṇe} again
\end{itemize}

\begin{tcolorbox}[breakable,colback=teal!4,colframe=teal!35!black,title={Before you move on}]
What does \mr{बंद} \mr{पडणे} mean? (\textquotedblleft{}To stop working; to break down\textquotedblright{}.) Say it once more.
\end{tcolorbox}
